\subsection{AFFINE LINEAR MAPPINGS}

DEFINITION 3. Given two affine spaces E, E' attached to two vector spaces T, T' over the same field K, a mapping u of E into E' is called an affine linear mapping (or an affine mapping) if, for every family \((x_{t})_{t\in I}\) of points of E and every family \((\lambda_{t})_{t\in I}\) such that \(\sum_{i\in I}\lambda_{i}=1\) ,

\[
u \left(\sum_ {i \in I} \lambda_ {i} x _ {i}\right) = \sum_ {i \in I} \lambda_ {i} u (x _ {i}).\tag{3}
\]

PROPOSITION 6. Let \(u\) be an affine mapping of \(E\) into \(E'\) . There exists one and only one linear mapping \(v\) of \(T\) into \(T'\) such that

\[
u (x + \mathbf {t}) = u (x) + v (\mathbf {t})
\]

for all \(x \in \mathbf{E}\) , \(\mathbf{t} \in \mathbf{T}\) .

Let \(a\) be any point of E. The mapping

\[
\mathbf {t} \mapsto u (a + \mathbf {t}) - u (a)
\]

is a linear mapping \(v_{a}\) of T into T', for we may write

\[
\begin{array}{c} a + \lambda \mathbf {t} = \lambda (a + \mathbf {t}) + (1 - \lambda) a \\ a + \mathbf {s} + \mathbf {t} = (a + \mathbf {s}) + (a + \mathbf {t}) - a \end{array}
\]

and it follows from (3) that \(v_{a}(\lambda \mathbf{t}) = \lambda v_{a}(\mathbf{t})\) and \(v_{a}(\mathbf{s} + \mathbf{t}) = v_{a}(\mathbf{s}) + v_{a}(\mathbf{t})\) . Moreover, if \(b\) is another point of \(\mathbf{E}\) , then \(v_{a} = v_{b}\) ; for the relation

\[
(a + \mathbf {t}) - a + b = b + \mathbf {t}
\]

implies

\[
u (a + \mathbf {t}) - u (a) + u (b) = u (b + \mathbf {t})
\]

that is \(u(a + \mathbf{t}) - u(a) = u(b + \mathbf{t}) - u(b)\) . Whence the existence of \(v\) ; the uniqueness is immediate.

v is called the linear mapping of T into T' associated with u. Conversely, for every linear mapping v of T into T' and every ordered pair of points \(a \in E\) , \(a' \in E'\) , it is immediately verified that

\[
x \mapsto a ^ {\prime} + v (x - a)
\]

is an affine mapping of E into \(E'\) whose associated linear mapping is v. To say that u is an affine mapping of E into \(E'\) therefore also means that, if an arbitrary point a in E and the point \(u(a)\) in \(E'\) are taken as origins, u is a linear mapping for the two vector spaces thus obtained.

Let \(\mathbf{E}''\) be a third affine space, \(\mathbf{T}''\) its translation space, \(u'\) an affine mapping of \(\mathbf{E}'\) into \(\mathbf{E}''\) and \(v'\) the linear mapping of \(\mathbf{T}'\) into \(\mathbf{T}''\) associated with \(u'\) . Clearly \(u' \circ u\) is an affine mapping of \(\mathbf{E}\) into \(\mathbf{E}''\) ; moreover, for \(a \in \mathbf{E}\) and \(\mathbf{t} \in \mathbf{T}\) ,

\[
u ^ {\prime} (u (a + \mathbf {t})) = u ^ {\prime} (u (a) + v (\mathbf {t})) = u ^ {\prime} (u (a)) + v ^ {\prime} (v (\mathbf {t}))
\]

and hence \(v'\circ v\) is the linear mapping of T into \(T''\) associated with \(u'\circ u\) . For an affine mapping u to be bijective, it is necessary and sufficient that the associated linear mapping v be so, and \(u^{-1}\) is then an affine mapping whose associated linear mapping is \(v^{-1}\) .

In particular, the affine bijections of E onto itself form a group G, called the affine group of E. The mapping which associates with \(u \in G\) the linear mapping v associated with u is, by the above, a homomorphism of G onto the linear group \(\mathbf{GL}(\mathrm{T})\) . If u is a translation, v is the identity and conversely. Hence, the kernel of the above homomorphism is the translation group T of E which is therefore a normal subgroup of G.

If \(u \in G\) , the automorphism \(t \mapsto utu^{-1}\) of T (where t is identified with the translation \(x \mapsto x + t\) ) is the linear mapping v associated with u. For \(x \in E\) and \(t \in T\) , by definition

\[
x + u \mathbf {t} u ^ {- 1} = u (u ^ {- 1} (x) + \mathbf {t}) = u (u ^ {- 1} (x)) + v (\mathbf {t}) = x + v (\mathbf {t})
\]

and hence \(utu^{-1} = v(t)\) .

Let \(a \in E\) and \(G_a\) be the subgroup of \(G\) consisting of the \(u \in G\) such that \(u(a) = a\) . If \(E\) is identified with \(T\) by taking \(a\) as origin, \(G_a\) is identified with \(\mathbf{GL}(T)\) . Every \(u \in G\) can be expressed uniquely in the form \(u = t_1 u_1\) (resp. in the form \(u = u_2 t_2\) ), where \(u_1, u_2\) are in \(G_a\) and \(t_1, t_2\) in \(T\) : for, writing \(t_1 = u(a) - a\) , \(u^{-1} t_1 \in G_a\) , whence the existence of \(u_1\) and \(t_1\) ; the existence of \(u_2\) and \(t_2\) is obtained analogously. The uniqueness follows from the fact that \(G_a \cap T\) reduces to the identity element of \(G\) . Moreover

\[
\mathbf {t} _ {1} u _ {1} = u _ {1} (u _ {1} ^ {- 1} \mathbf {t} _ {1} u _ {1})
\]

whence \(u_{2}=u_{1}\) , \(t_{2}=u_{1}^{-1}t_{1}u_{1}\) . Finally, the linear mappings associated with u and \(u_{1}\) are the same and hence, if as above \(G_{a}\) is identified with \(\mathbf{GL}(T)\) , \(u_{1}\) is the linear mapping from T to itself associated with u. It is thus seen that G is the semi-direct product of \(G_{a}\) by T (I, §6, no. 1).

Let E, E' be two affine spaces over K. The direct (resp. inverse) image of a linear variety of E (resp. E') under an affine mapping u of E into E' is a linear variety of E' (resp. E); the rank of u is by definition the dimension of \(u(\mathrm{E})\) ; it is equal to the rank of the linear mapping associated with u. If V, V' are linear varieties of the same finite dimension m in E, E' respectively, there exists an affine mapping \(u\) of \(\mathbf{E}\) into \(\mathbf{E}'\) such that \(u(\mathbf{V}) = \mathbf{V}'\) : taking as origins in \(\mathbf{E}\) and \(\mathbf{E}'\) points of \(\mathbf{V}\) and \(\mathbf{V}'\) respectively, then taking in \(\mathbf{E}\) (resp. \(\mathbf{E}'\) ) a basis whose first \(m\) vectors form a basis of \(\mathbf{V}\) (resp. \(\mathbf{V}'\) ), the proposition follows immediately from § 1, no. 11, Corollary 3 to Proposition 17.

As the field K has canonically a left vector space structure (of dimension 1) over K, it can be considered as an affine space of dimension 1. An affine mapping of an affine space D (over K) into the affine space K is also called an affine linear function (or an affine function). If a point \(a\) is taken as origin in E, every affine function on E can then be written uniquely as \(x \mapsto \alpha + v(x)\) , where \(\alpha \in \mathbf{K}\) and \(v\) is a linear form on the vector space E thus obtained; the affine functions on E therefore form a right vector space over K of dimension \(1 + \dim \mathbf{E}\) . If \(u\) is a non-constant affine function on E and \(\lambda \in \mathbf{K}\) , the set of \(x \in \mathbf{E}\) satisfying the equation \(u(x) = \lambda\) is a hyperplane; conversely, for every

hyperplane H in E, there exists an affine function \(u_{0}\) on E such that \(\mathbf{H} = \bar{u}_{0}^{-1}(0)\) and every affine function u such that \(\mathbf{H} = \bar{u}^{-1}(0)\) is of the form \(u_{0}\mu\) , where \(\mu \in K\) (§7, no. 5, Proposition 11). If u is an affine function on E, the hyperplanes with equations \(u(x) = \alpha\) and \(u(x) = \beta\) are parallel.

